\section{Implications}
\label{sec:ch9-implications}

The answer to \hyperref[rq:ch8-rq1]{\textbf{RQ1}} in Section~\ref{subsec:ch8-rq1-answer} shows that a probe on the hidden states of a coordinator selects the correct sub-agent more reliably than the coordinator itself. The probe is a linear classifier on the hidden state at the last prompt token, which the coordinator already computes, so a multi-agent system obtains a more accurate routing decision without an additional model call. The gain grows as the correct sub-agent becomes more difficult to determine from the wording of the request, which is where the coordinator misroutes most often. The gain is also smaller when the request is part of a multi-turn conversation.

The answer to \hyperref[rq:ch8-rq2]{\textbf{RQ2}} in Section~\ref{subsec:ch8-rq2-answer} shows that the same gain is not obtained by the three text-only classifiers of this work, which are logistic regressions trained on a bag-of-words, a TF-IDF and an embedding representation of the request text and need no access to the hidden states. As the correct sub-agent becomes more difficult to determine from the wording of the request, the accuracy of each text-only classifier falls more than the accuracy of the probe. On these requests, the hidden states of the coordinator therefore carry routing information that the representations of the request text do not provide, and this information is what justifies reading the hidden states. As the difficulty of the requests increases, the probe is the classifier that retains the most accuracy to evaluate the routing decision of the coordinator.

The answer to \hyperref[rq:ch8-rq3]{\textbf{RQ3}} in Section~\ref{subsec:ch8-rq3-answer} shows that the optimal layer for the probe falls at a consistent relative depth across the models of this work. For a new coordinator model, it is therefore likely to fall at a similar depth, and the layer sweep aimed at finding it can be restricted around that depth.

Because the probe is more accurate than the coordinator, it can be used as an evaluation flag on the routing decision of the coordinator. The answer to \hyperref[rq:ch8-rq4]{\textbf{RQ4}} in Section~\ref{subsec:ch8-rq4-answer} shows that the disagreement flag, raised when the probe and the coordinator select different sub-agents, marks the requests that the coordinator is likely to misroute, with a precision above the base error rate of the coordinator in every run. When the flag is raised, the execution of the request can be stopped (Section~\ref{sec:early-exit-feedback-loop}), or the request can be kept for later review by a human. On the single-turn dataset, the flag is raised on between 56.0\% and 79.7\% of the misroutes of the coordinator, depending on the model, and between 68.1\% and 82.6\% of the flagged requests are misroutes of the coordinator. On the multi-turn dataset, the flag is raised on between 54.0\% and 77.8\% of the misroutes, and between 59.8\% and 80.5\% of the flagged requests are misroutes of the coordinator.

Because the probe is more accurate than the coordinator, the request could also be handed off to the sub-agent that the probe predicts, instead of the sub-agent that the coordinator selects. The probe would then replace the routing decision of the coordinator, instead of serving as an evaluation flag on it. With this replacement, a smaller coordinator can give a more accurate routing decision than a larger one: on both datasets, the probes on Qwen3 4B Instruct 2507 and on Gemma 4 E2B IT are more accurate than the Llama 3.3 70B Instruct and the Qwen2.5 72B Instruct coordinators (Tables~\ref{tab:ch7-qwen3-4b}, \ref{tab:ch7-gemma4-e2b-it}, \ref{tab:ch7-llama33-70b} and~\ref{tab:ch7-qwen25-72b}). Ratnakar et al. reach the same conclusion with LatentGate, where a linear probe on the hidden states of Qwen2.5 1.5B Instruct routes queries more accurately than prompt-based routers built on GPT-4.1, GPT-4.1-nano and Gemini 2.5 Flash \cite{ratnakarLatentGateLowLatencySemantic2026}. However, this replacement may not be possible when the coordinator has responsibilities other than routing, such as conducting the conversation with the user, producing the final response from the response of the selected sub-agent (Section~\ref{sec:coordinator}), or applying guardrails (Section~\ref{sec:security-evaluation}).
